\documentclass[11pt,a4paper]{article}

% ---------------------------------------------------------------------------
% Rapport de déploiement — Frontend GANDAL
% Compilation : pdflatex rapport_deploiement_frontend.tex (x2 pour la TOC)
% ---------------------------------------------------------------------------

\usepackage[utf8]{inputenc}
\usepackage[T1]{fontenc}
\usepackage[french]{babel}
\usepackage{lmodern}
\usepackage[a4paper,margin=2.5cm]{geometry}
\usepackage{graphicx}
\usepackage{xcolor}
\usepackage{titlesec}
\usepackage{fancyhdr}
\usepackage{booktabs}
\usepackage{tabularx}
\usepackage{enumitem}
\usepackage{listings}
\usepackage{hyperref}

% --- Couleurs ---
\definecolor{gandalblue}{RGB}{20,40,90}
\definecolor{gandalaccent}{RGB}{0,120,160}
\definecolor{codebg}{RGB}{245,245,247}
\definecolor{codekw}{RGB}{160,30,90}
\definecolor{codecmt}{RGB}{100,110,120}
\definecolor{okgreen}{RGB}{30,130,60}

\hypersetup{
  colorlinks=true,
  linkcolor=gandalblue,
  urlcolor=gandalaccent,
  pdftitle={Rapport de deploiement - Frontend GANDAL},
  pdfauthor={Equipe DevOps GANDAL}
}

% --- En-têtes / pieds de page ---
\pagestyle{fancy}
\fancyhf{}
\fancyhead[L]{\small\textsc{GANDAL — Déploiement Frontend}}
\fancyhead[R]{\small\thepage}
\fancyfoot[C]{\small ENSPY — Département de Génie Informatique}
\renewcommand{\headrulewidth}{0.4pt}
\renewcommand{\footrulewidth}{0.2pt}

% --- Titres de section ---
\titleformat{\section}{\Large\bfseries\color{gandalblue}}{\thesection}{0.6em}{}
\titleformat{\subsection}{\large\bfseries\color{gandalaccent}}{\thesubsection}{0.6em}{}

% --- Style du code ---
\lstdefinestyle{shell}{
  backgroundcolor=\color{codebg},
  basicstyle=\ttfamily\footnotesize,
  keywordstyle=\color{codekw},
  commentstyle=\color{codecmt}\itshape,
  breaklines=true,
  frame=single,
  rulecolor=\color{gray!40},
  framesep=6pt,
  xleftmargin=8pt,
  showstringspaces=false,
  columns=fullflexible,
}
\lstset{style=shell}

% --- Macro encadré statut ---
\newcommand{\statut}[1]{\textcolor{okgreen}{\textbf{#1}}}

\begin{document}

% =========================== PAGE DE TITRE ===========================
\begin{titlepage}
\centering
\vspace*{2cm}
{\color{gandalblue}\rule{\linewidth}{1.5pt}}\\[0.6cm]
{\Huge\bfseries\color{gandalblue} Rapport de Déploiement}\\[0.4cm]
{\LARGE Plateforme GANDAL — Frontend Web}\\[0.3cm]
{\large\itshape Mise en production sur le data center local}\\[0.4cm]
{\color{gandalblue}\rule{\linewidth}{1.5pt}}\\[1.5cm]

\begin{tabular}{rl}
\textbf{Application}     & Frontend Next.js 16 (App Router)\\[4pt]
\textbf{Serveur cible}   & \texttt{root@10.50.30.102} (Debian 12)\\[4pt]
\textbf{URL publique}    & \url{http://gandal-app.10.50.30.102.nip.io}\\[4pt]
\textbf{Méthode}         & Conteneurisation Docker + reverse proxy nginx\\[4pt]
\textbf{Date}            & 11 juin 2026\\[4pt]
\textbf{Statut}          & \statut{Déployé et vérifié}\\
\end{tabular}

\vfill
{\large\textbf{École Nationale Supérieure Polytechnique de Yaoundé (ENSPY)}}\\[0.2cm]
{Département de Génie Informatique — Promotion GI27}\\[0.2cm]
{Université de Yaoundé I}\\[1cm]
{\small \copyright\ 2026 GANDAL Data Center}
\end{titlepage}

% =========================== SOMMAIRE ===========================
\tableofcontents
\thispagestyle{fancy}
\newpage

% =========================== 1. CONTEXTE ===========================
\section{Contexte et objectif}

GANDAL est une plateforme académique de gestion centralisée de projets étudiants,
adossée à un data center local du Département de Génie Informatique. Le présent
document décrit la mise en production de la \textbf{couche frontend} (interface web)
sur le serveur d'hébergement \texttt{10.50.30.102}, aux côtés de l'API backend déjà
en service.

\subsection{Périmètre}
\begin{itemize}[leftmargin=1.4em]
  \item Construction d'une image Docker de l'application Next.js.
  \item Démarrage du conteneur sur l'infrastructure existante.
  \item Exposition via le reverse proxy nginx mutualisé.
  \item Configuration du partage de ressources cross-origin (CORS) côté API.
  \item Vérification fonctionnelle de bout en bout.
\end{itemize}

% =========================== 2. ENVIRONNEMENT ===========================
\section{Environnement cible}

\begin{table}[h]
\centering
\renewcommand{\arraystretch}{1.3}
\begin{tabularx}{\linewidth}{@{}lX@{}}
\toprule
\textbf{Élément} & \textbf{Valeur} \\
\midrule
Hôte & \texttt{omega-test-3001}, Debian GNU/Linux 12, x86\_64 \\
Moteur de conteneurs & Docker 29.5.3 + Docker Compose v5.1.4 \\
Node.js sur l'hôte & \textit{Absent} (d'où le choix de la conteneurisation) \\
Reverse proxy & Conteneur \texttt{iwm-gateway} (nginx 1.29.8-alpine), port 80 \\
Backend existant & \texttt{gandal-api} (port 8000) + \texttt{gandal-db} (PostgreSQL 17) \\
Réseau Docker partagé & \texttt{iwm-vm-network} (résolution par nom de service) \\
\bottomrule
\end{tabularx}
\caption{Caractéristiques du serveur de production.}
\end{table}

\subsection{Justification du choix Docker}
L'ensemble de l'infrastructure du serveur est conteneurisé : le reverse proxy nginx
route le trafic vers des \emph{noms de service Docker} sur le réseau interne. Le serveur
ne disposant pas de runtime Node.js, et un export statique étant impossible (l'application
utilise un \emph{middleware} d'authentification exécuté côté serveur), la conteneurisation
est la solution cohérente et la moins intrusive, alignée sur le déploiement du backend.

% =========================== 3. ARCHITECTURE ===========================
\section{Architecture de déploiement}

Le frontend est exposé sur un \textbf{hôte virtuel dédié} distinct de l'API, le routage
étant assuré par la directive \texttt{map \$host \$svc} du reverse proxy.

\begin{center}
\footnotesize
\begin{tabular}{c}
\fbox{\parbox{0.92\linewidth}{\centering
\textbf{Navigateur}\\[2pt]
$\downarrow$ \texttt{http://gandal-app.10.50.30.102.nip.io} \quad / \quad \texttt{http://gandal.10.50.30.102.nip.io}\\[6pt]
\textbf{iwm-gateway} (nginx, :80) — routage par nom d'hôte\\[6pt]
$\swarrow$ \texttt{gandal-app.*} \hspace{3.5cm} $\searrow$ \texttt{gandal.*}\\[6pt]
\fbox{\texttt{gandal-dev-frontend:3000}} \hspace{1.2cm} \fbox{\texttt{gandal-api:8000}}\\[4pt]
(Next.js / \texttt{next start}) \hspace{2.0cm} (FastAPI) $\rightarrow$ \texttt{gandal-db:5432}
}}
\end{tabular}
\end{center}

\begin{itemize}[leftmargin=1.4em]
  \item \texttt{gandal-app.10.50.30.102.nip.io} $\rightarrow$ conteneur frontend (port 3000).
  \item \texttt{gandal.10.50.30.102.nip.io} $\rightarrow$ API backend (port 8000), inchangé.
  \item Les appels API du navigateur étant cross-origin, le backend autorise
        explicitement l'origine du frontend (cf. \S\ref{sec:cors}).
\end{itemize}

% =========================== 4. ARTEFACTS ===========================
\section{Artefacts de déploiement}

\subsection{Dockerfile (multi-stage)}
Image construite en deux étapes : compilation puis exécution avec l'utilisateur
non privilégié \texttt{node}. L'approche s'appuie sur les scripts du projet
(\texttt{next build} / \texttt{next start}), robuste vis-à-vis des spécificités de
cette distribution de Next.js.

\begin{lstlisting}[language=bash]
FROM node:20-alpine AS build
WORKDIR /app
COPY package.json package-lock.json ./
RUN npm ci
COPY . .
ENV NEXT_TELEMETRY_DISABLED=1
RUN npm run build
RUN npm prune --omit=dev

FROM node:20-alpine AS runner
WORKDIR /app
ENV NODE_ENV=production PORT=3000 HOSTNAME=0.0.0.0
COPY --from=build --chown=node:node /app/package.json ./package.json
COPY --from=build --chown=node:node /app/node_modules ./node_modules
COPY --from=build --chown=node:node /app/.next ./.next
COPY --from=build --chown=node:node /app/public ./public
COPY --from=build --chown=node:node /app/next.config.ts ./next.config.ts
USER node
EXPOSE 3000
CMD ["npm", "run", "start"]
\end{lstlisting}

\subsection{Orchestration (docker-compose.yml)}
\begin{lstlisting}[language=bash]
name: gandal-frontend
services:
  frontend:
    build: { context: ., dockerfile: Dockerfile }
    image: gandal-dev-frontend:latest
    container_name: gandal-dev-frontend
    restart: unless-stopped
    environment: { NODE_ENV: production, PORT: "3000", HOSTNAME: 0.0.0.0 }
    expose: ["3000"]
    networks: [gateway-net]
networks:
  gateway-net:
    external: true
    name: iwm-vm-network
\end{lstlisting}

% =========================== 5. PROCEDURE ===========================
\section{Procédure de déploiement}

\subsection{Transfert des sources}
Les sources (\(\approx\)11 Mo, assets statiques inclus) sont synchronisées vers
\texttt{/root/gandal-dev-frontend} :
\begin{lstlisting}[language=bash]
rsync -az --delete --exclude=node_modules --exclude=.next --exclude=.git \
      ./ root@10.50.30.102:/root/gandal-dev-frontend/
\end{lstlisting}

\subsection{Construction et démarrage}
La construction est réalisée directement sur le serveur (build isolé dans le
conteneur, indépendant de l'environnement hôte) :
\begin{lstlisting}[language=bash]
cd /root/gandal-dev-frontend && docker compose up -d --build
\end{lstlisting}
Résultat : image \texttt{gandal-dev-frontend:latest} (\(\approx\)911 Mo), conteneur
\texttt{gandal-dev-frontend} en état \texttt{running} sur le réseau \texttt{iwm-vm-network}.

\subsection{Routage nginx}
Ajout d'une entrée dans \texttt{/root/iwm/ops/nginx/gateway.conf} :
\begin{lstlisting}[language=bash]
map $host $svc {
    default               "";
    "~^gandal-app\."      "gandal-dev-frontend:3000";   # <-- ajout
    "~^gandal(-api)?\."   "gandal-api:8000";
    ...
}
\end{lstlisting}

\subsection{Configuration CORS}
\label{sec:cors}
L'origine du frontend est ajoutée à \texttt{CORS\_ORIGINS} dans
\texttt{/gandal-dev/.env}, puis le service API est recréé pour prise en compte :
\begin{lstlisting}[language=bash]
CORS_ORIGINS=...,http://gandal-app.10.50.30.102.nip.io
cd /gandal-dev && docker compose -f docker-compose.prod.yml up -d api
\end{lstlisting}

% =========================== 6. INCIDENT ===========================
\section{Incident rencontré et résolution}

\subsection{Symptôme}
Après ajout de la règle de routage, l'URL du frontend renvoyait \texttt{HTTP 404 —
Unknown host}, alors que la configuration sur disque était correcte et que la
résolution DNS interne du conteneur fonctionnait.

\subsection{Cause racine}
Le conteneur \texttt{iwm-gateway} monte un \emph{fichier unique}
(\texttt{gateway.conf}) en \emph{bind mount}. Une modification via \texttt{sed -i}
remplace l'\emph{inode} du fichier ; le conteneur reste alors accroché à l'ancien
inode. Les rechargements \texttt{nginx -s reload} relisaient donc une configuration
figée, sans la nouvelle règle.

\subsection{Résolution}
Validation de la configuration dans un conteneur nginx jetable, puis redémarrage du
gateway pour re-lier le fichier :
\begin{lstlisting}[language=bash]
docker run --rm -v /root/iwm/ops/nginx/gateway.conf:\
/etc/nginx/conf.d/default.conf:ro nginx:1.29.8-alpine nginx -t   # OK
docker restart iwm-gateway
\end{lstlisting}

\subsection{Recommandation}
Pour toute modification ultérieure du gateway, privilégier une réécriture
\emph{in-place} préservant l'inode, ou prévoir le redémarrage du conteneur. À terme,
monter un \emph{répertoire} plutôt qu'un fichier unique élimine ce risque.

% =========================== 7. VERIFICATION ===========================
\section{Vérification fonctionnelle}

\begin{table}[h]
\centering
\renewcommand{\arraystretch}{1.35}
\begin{tabularx}{\linewidth}{@{}Xll@{}}
\toprule
\textbf{Contrôle} & \textbf{Attendu} & \textbf{Résultat} \\
\midrule
Page d'accueil via gateway & HTTP 200 & \statut{200} (HTML, titre GANDAL) \\
Asset statique \texttt{/logo.png} & HTTP 200 & \statut{200} (112 Ko) \\
API \texttt{gandal.*/health} & HTTP 200 & \statut{200} (backend intact) \\
Préflight CORS depuis le front & En-tête autorisé & \statut{allow-origin présent} \\
Persistance conteneur & \texttt{restart=unless-stopped} & \statut{Actif} \\
\bottomrule
\end{tabularx}
\caption{Synthèse des tests post-déploiement.}
\end{table}

\noindent En-tête CORS retourné par l'API pour l'origine du frontend :
\begin{lstlisting}[language=bash]
access-control-allow-origin: http://gandal-app.10.50.30.102.nip.io
access-control-allow-credentials: true
access-control-allow-methods: DELETE, GET, HEAD, OPTIONS, PATCH, POST, PUT
\end{lstlisting}

% =========================== 8. EXPLOITATION ===========================
\section{Exploitation et maintenance}

\subsection{Commandes usuelles}
\begin{lstlisting}[language=bash]
# Etat / logs
docker ps --filter name=gandal-dev-frontend
docker logs -f gandal-dev-frontend

# Mise a jour (apres rsync des nouvelles sources)
cd /root/gandal-dev-frontend && docker compose up -d --build

# Redemarrage
docker restart gandal-dev-frontend
\end{lstlisting}

\subsection{Procédure de rollback}
\begin{enumerate}[leftmargin=1.6em]
  \item Restaurer la sauvegarde du gateway : \texttt{gateway.conf.bak.*} puis
        \texttt{docker restart iwm-gateway}.
  \item Restaurer \texttt{/gandal-dev/.env.bak.*} et recréer l'API.
  \item Arrêter le frontend : \texttt{docker compose down} dans
        \texttt{/root/gandal-dev-frontend}.
\end{enumerate}

% =========================== 9. CONCLUSION ===========================
\section{Conclusion}

Le frontend GANDAL est déployé en production sur \texttt{10.50.30.102} et accessible
à l'adresse \url{http://gandal-app.10.50.30.102.nip.io}. L'intégration au reverse
proxy mutualisé et la configuration CORS permettent une communication complète avec
l'API existante, sans régression sur les autres services hébergés. L'ensemble est
géré par Docker Compose, avec redémarrage automatique et procédure de rollback
documentée.

\subsection*{Points d'amélioration suggérés}
\begin{itemize}[leftmargin=1.4em]
  \item Mise en place de TLS (HTTPS) sur le gateway.
  \item Intégration continue : build et publication de l'image vers un registre
        (ex. GHCR) pour réduire la charge de build sur le serveur.
  \item Surveillance de la mémoire (RAM limitée, absence de swap sur l'hôte).
\end{itemize}

\vfill
\begin{center}
\small\itshape Document généré dans le cadre du déploiement du 11 juin 2026.
\end{center}

\end{document}
